\section{Operational concerns and diagnostic scope}
\par\smallskip\noindent\begin{tabularx}{\linewidth}{@{}p{.27\linewidth}XX@{}}\toprule
Concern & Relevant observation & Interpretation limit\\\midrule
Gearbox vibration & RMS velocity in mm/s; trend with load & No spectral or bearing-frequency measurement\\
Bearing deterioration & Vibration increase; noise reported by operator & No measured defect frequencies or inspection labels\\
Oil overheating & Oil temperature, load and vibration & No viscosity, debris or oil chemistry signal\\
Excessive screw torque & Torque percentage, speed, feed, pressure & Recipe resistance and obstruction can look similar\\
Barrel instability & Aggregate barrel temperature & No zone-wise heater/cooler diagnosis\\
Melt-pressure abnormality & Pressure, melt temperature, throughput & No die geometry or material viscosity measurement\\
Unstable feeding & Feed-rate change with torque variation & No direct hopper/feeder blockage sensor\\
Excessive motor current & Current, power and torque & No phase currents or electrical fault waveform\\\bottomrule
\end{tabularx}\par\smallskip
These are engineering hypotheses, not validated procedures or machine-specific thresholds. The live RAG issues contain the same eight concern categories, with symptom and possible-cause text. Those records are illustrative maintenance knowledge: they are not verified incidents, work orders or manufacturer fault diagnoses. The database's \texttt{possible\_cause} field deliberately preserves this distinction.
\section{Four different questions}
\textbf{Fault detection} asks whether observed behavior appears abnormal. The binary synthetic classifier addresses this limited question. \textbf{Fault diagnosis} identifies a physical cause and often requires discriminating tests, inspection and domain expertise. \textbf{Predictive maintenance} forecasts deterioration or useful life sufficiently to plan intervention. The implementation has no remaining-useful-life model or forecast horizon. \textbf{Maintenance decision support} organizes measurements, model outputs, documents and uncertainty for a human decision. These concepts follow the separation between diagnostics and prognostics in condition-based maintenance research~\cite{jardine}.
\section{Requirements derived from the operator's question}
The assistant should identify the correct asset, state when each observation was measured, preserve recipe and units, report the actual model version, retrieve relevant passages, and distinguish evidence from interpretation. A recommendation should say which inspection merits consideration and what is still unknown. It should not turn a binary label into an automatic shutdown, repair instruction or certification that operation is safe.

Historical behavior is useful only when comparable operating conditions are identified. A current reading at high load should not be compared indiscriminately with a low-load recipe baseline. The repository does not contain a complete physical asset record, screw configuration, rated motor power, die configuration, approved alarm envelope or plant safety case. Its sensor validation bounds are API input limits, not Coperion operating limits.
\verdict{Combining modalities can improve the explanation available to an operator, but correct diagnosis still requires independently verified asset identity, measurement validity and qualified maintenance judgment.}
\source{\path{rag-server/rag/services.py}: issue CRUD and serialization; \path{zsk-service/zsk/zsk_model.py}: \texttt{LABELS}; live issue catalog (eight records).}
